\emph{Premise.} Hypothesis (P) of the statement is Premise \claimref{synthesis.premise-positivity-of-values}; every statement below about the values of states is under it, relative to the fixed split of the generators.

\emph{Preliminaries.} The adjoint, observables, configurations and (P) are those of Theorem \claimref{synthesis.chsh-bound-for-two-subsystems-of-a-clifford-algebra-by-a-rational-certificate}, whose preliminaries give: the blades are a basis, $\langle xy^{\dagger}\rangle_{0}$ is the sum of products of coefficients, $\langle ab\rangle_{0} = \langle ba\rangle_{0}$, a self-adjoint blade has square $1$, and $\omega_{\rho}(x) = \langle\rho x\rangle_{0}/\langle\rho\rangle_{0}$ is a state for $\rho = \sum y_{r}y_{r}^{\dagger} \ne 0$. Its clause (i) is $C^{2} = 4 - [A_{1},A_{2}][B_{1},B_{2}]$ and its clause (iii) gives $\omega(h)^{2} \le \omega(h^{2})$ for self-adjoint $h$ under (P). \emph{The twist.} Each $p_{j} = f_{j}e_{j}$ is self-adjoint of square $1$, anticommutes with $e_{j}$ and $f_{j}$ and commutes with the other generators; so $\Gamma_{A}$ has square $1$, anticommutes with the generators of $\Cl(A)$, commutes with those of $\Cl(B)$, and for generators $y, y'$ of $\Cl(B)$ and $x$ of $\Cl(A)$: $(y\Gamma_{A})(y'\Gamma_{A}) = yy'$ and $x\,y\Gamma_{A} = -yx\Gamma_{A} = y\Gamma_{A}x$. These are the relations of the twisted copy of Theorem \claimref{tower.commuting-copies} (there with ground algebra $\Cl(B)$ and the planes of $A$), proved here in three lines.

\emph{(i).} Distinct generators anticommute, so homogeneous $x \in \Cl(A)$, $y \in \Cl(B)$ satisfy $xy = (-1)^{|x||y|}yx$ and even elements commute. By the twist relations a product of an even number of twisted generators is the product of the untwisted ones, so the even part of $B'$ is the even part of $\Cl(B)$. $\Gamma$ anticommutes with every generator, so an element of $\Cl(A)$ or $\Cl(B)$ commutes with $\Gamma$ exactly when it is even.

\emph{(ii).} The blades of the algebra of mode $j$ are $1$, $e_{j}$, $f_{j}$, $p_{j}$; $f_{j}^{\dagger} = -f_{j}$ and $p_{j}^{\dagger} = e_{j}^{\dagger}f_{j}^{\dagger} = -e_{j}f_{j} = p_{j}$, so the self-adjoint elements are $s + tp_{j} + ue_{j}$. Since $p_{j}$ and $e_{j}$ anticommute and have square $1$, $(s + tp_{j} + ue_{j})^{2} = s^{2} + t^{2} + u^{2} + 2s(tp_{j} + ue_{j})$, which is $1$ exactly when $t = u = 0$, $s = \pm 1$, or $s = 0$, $t^{2} + u^{2} = 1$. The element is even exactly when $u = 0$. In the twisted copy the same computation holds with $e_{j}\Gamma_{A}$ in place of $e_{j}$ (self-adjoint, of square $1$, anticommuting with $p_{j}$, and odd). So the even observables are $\pm 1$, $\pm p_{j}$; they commute, $[A_{1},A_{2}] = 0$, $C^{2} = 4$ and $\omega(C)^{2} \le 4$. Nothing was assumed about the other party beyond $A_{i}B_{j} = B_{j}A_{i}$.

\emph{(iii).} $e_{1}f_{2}$ is self-adjoint ($f_{2}^{\dagger}e_{1}^{\dagger} = -f_{2}e_{1} = e_{1}f_{2}$), so of square $1$, and it anticommutes with $p_{1} = f_{1}e_{1}$, which contains $e_{1}$ and not $f_{2}$; likewise for $B$. All four are products of two generators, hence even, and commute across the parties by (i). So they are a configuration of blades with anticommuting pairs, and the blade lemma applies: $K = -A_{1}A_{2}B_{1}B_{2} = -(p_{1}e_{1}f_{2})(p_{3}e_{3}f_{4}) = -f_{1}f_{2}f_{3}f_{4}$, and the functionals $\omega_{r}$ satisfy (P) for $8r^{2} \le 1$ with $\omega_{r}(C) = 8r$. The element $(1 + K)(1 + rC)$ is even, and $\langle\Gamma a\Gamma\rangle_{0} = \langle a\rangle_{0}$, so $\omega_{r}$ is even.

\emph{(iv).} By the twist relations $e_{2}p_{1}$ has square $1$ and commutes with $\Cl(A)$; it is self-adjoint since $e_{2}$ and $p_{1}$ are and commute; it anticommutes with $p_{2}$ as $e_{2}$ does; $p_{2}$ commutes with $\Cl(A)$. So the four elements are a configuration. Put $X = p_{1}p_{2}$ and $Y = -f_{1}e_{2}$: both are self-adjoint blades of square $1$, they commute, and $XY = -e_{1}f_{2}$; so $\rho = \tfrac14(1 + X)(1 + Y)$ is a self-adjoint idempotent, $\rho = \rho\rho^{\dagger}$, $4\langle\rho\rangle_{0} = 1$, and $\omega = 4\langle\rho\,\cdot\,\rangle_{0}$ is a state by the preliminaries; $\rho$ is even, so $\omega$ is even. The value of a blade $b$ of square $1$ is its coefficient in $4\rho$: it is $1$ for $p_{1}p_{2}$ and for $e_{1}\,e_{2}p_{1} = e_{2}f_{1} = -f_{1}e_{2}$, and $0$ for the odd blades $p_{1}\,e_{2}p_{1} = e_{2}$ and $e_{1}p_{2}$. With $B_{\pm} = zp_{2} \pm x\,e_{2}p_{1}$ (observables by (ii)): $\omega(C) = \omega(p_{1}(B_{+} + B_{-})) + \omega(e_{1}(B_{+} - B_{-})) = 2z\,\omega(p_{1}p_{2}) + 2x\,\omega(e_{1}e_{2}p_{1}) = 2(z + x)$, and $(z + x)^{2} = 1 + 2xz > 1$ for $x, z > 0$. \emph{General amplitudes.} With $v_{j} = \tfrac12(e_{j} + f_{j})$, $w_{j} = \tfrac12(e_{j} - f_{j})$ the stored pair state is $\tfrac14(1 + p_{1}p_{2})(1 + (c^{2} - s^{2})p_{1}) + cs(w_{1}w_{2} + v_{2}v_{1})$ and $w_{1}w_{2} + v_{2}v_{1} = -\tfrac12(e_{1}f_{2} + f_{1}e_{2})$, which is $\rho_{c,s}$; it equals $\tfrac14(1 + X)(1 + (c^{2} - s^{2})p_{1} + 2cs\,Y)$, a self-adjoint idempotent because $(c^{2} - s^{2})p_{1} + 2cs\,Y$ is self-adjoint of square $1$ ($p_{1}$ and $Y$ anticommute) and commutes with $X$. The coefficient of $f_{1}e_{2}$ in $4\rho_{c,s}$ is $-2cs$, so $\omega(e_{1}e_{2}p_{1}) = 2cs$, $\omega(p_{1}p_{2}) = 1$, the odd blades have value $0$, and $\omega(C) = 2(z + 2cs\,x)$. By the Cauchy-Schwarz inequality $(z + 2cs\,x)^{2} \le (1 + 4c^{2}s^{2})(x^{2} + z^{2})$, with equality when $(x, z)$ is proportional to $(2cs, 1)$; the supremum over the directions is $2\sqrt{1 + 4c^{2}s^{2}}$. \emph{Parity.} $\Gamma$ anticommutes with every generator, so with $e_{1}$ and with $e_{2}p_{1}$ (three generators), and the parity average $E(x) = \tfrac12(x + \Gamma x\Gamma)$, which is the average of Theorem \claimref{singularity.descent-parity-average} with the total parity $\Gamma$ in place of its single-mode $\sigma$, sends them to $0$; $E(p_{j}) = p_{j}$; the CHSH combination of the averaged elements is $2z\,p_{1}p_{2}$, of value $2z$. The last statement is (ii) applied to either party.

\emph{(v).} $e_{1}e_{2} = -e_{2}e_{1}$ and $(e_{1}e_{2})^{\dagger} = e_{2}e_{1} = -e_{1}e_{2}$, so the reality clause of (P) gives $\omega(e_{1}e_{2}) = -\omega(e_{1}e_{2}) = 0$. (Positivity alone does not give $\omega(x^{\dagger}) = \omega(x)$ over real coefficients, which is why the reality clause is part of (P).) Put $X' = e_{1}p_{2}$, $Y' = p_{1}e_{2}$, $Z' = p_{1}p_{2}$: each is a product of two commuting self-adjoint elements of square $1$, so is self-adjoint of square $1$, and $e_{1}e_{2}$ is the only term of the untwisted element that is not self-adjoint, so its self-adjoint part is $H = X' + Y' - Z'$. $X'Z' = e_{1}p_{1} = -p_{1}e_{1} = -Z'X'$, likewise $Y'Z' = -Z'Y'$, and $X'Y' = e_{1}p_{1}\,p_{2}e_{2}$, $Y'X' = p_{1}e_{2}e_{1}p_{2} = -p_{1}e_{1}\,e_{2}p_{2} = -X'Y'$. So $H^{2} = 3$ and $\omega(H)^{2} \le \omega(H^{2}) = 3$.
